
% =========================================================
% CHAPTER 2: REQUIREMENTS SPECIFICATION
% =========================================================
\chapter{Requirements Analysis}

\section{Purpose}
The purpose of the Smart Transaction Sorter and Analytics System is to provide users with an intelligent, automated way to analyze their personal financial spending. This web application uses Artificial Intelligence (AI) to categorize raw transaction data (from CSV files) and instantly convert it into clear, visual charts, eliminating the need for hours of frustrating manual work.

\section{Project Scope}
The scope of this project is a standalone, web-based tool that accepts a financial history file (CSV), uses AI to categorize every line based on user-defined labels, and presents the results on a BI Dashboard. The system includes User authentication, Category Management, CSV Upload/Validation, AI Sorting Logic, Data Persistence, and a Visualization Dashboard.

\begin{itemize}
	\item \textbf{In Scope:} User login, Category creation, CSV parsing, AI classification, Dashboard charts.
	\item \textbf{Out of Scope:} Direct integration (API) with bank/wallet accounts, payment processing, or real-time transaction tracking.
\end{itemize}

\section{Definitions, Acronyms, and Abbreviations}
\renewcommand{\arraystretch}{1.5}
\begin{longtable}{|p{4cm}|p{10.5cm}|}
	\hline
	\textbf{Term/Abbreviation} & \textbf{Definition} \\
	\hline
	\endfirsthead
	\hline
	\textbf{Term/Abbreviation} & \textbf{Definition} \\
	\hline
	\endhead
	STS & Smart Transaction Sorter (The name of this application). \\ \hline
	AI & Artificial Intelligence (The component that performs automatic classification). \\ \hline
	CSV & Comma-Separated Values (The standard file format for data input/upload). \\ \hline
	BI & Business Intelligence (Refers to the analytical dashboard and visualization tools). \\ \hline
	DB & Database (The system storage where user and financial data are saved). \\ \hline
\end{longtable}

\section{Document Overview}
This document describes the requirements for the Smart Transaction Sorter and Analytics System (STS). It serves as the single reference point for what the system must do and how well it must perform. The STS acts as a complimentary Business Intelligence layer for existing digital wallets and bank applications that currently only offer raw data exports.

\section{Functional Requirements}

\subsection{Security Management}

\subsubsection{Process Sign Up}
\begin{longtable}{|p{2.5cm}|p{12cm}|}
	\hline
	\textbf{ID} & \textbf{Requirement Description} \\
	\hline
	SRS-1 & The system shall allow new users to securely register by providing a username, email, and password. \\ \hline
	SRS-2 & The system must validate that the email address is unique and not already registered. \\ \hline
	SRS-3 & The system shall store user passwords using a strong one-way hash (never in plain text). \\ \hline
	SRS-4 & Upon successful registration, the system shall create a dedicated, isolated database scope for the user. \\ \hline
\end{longtable}

\subsubsection{Process Sign In / Logout}
\begin{longtable}{|p{2.5cm}|p{12cm}|}
	\hline
	\textbf{ID} & \textbf{Requirement Description} \\
	\hline
	SRS-5 & The system shall allow registered users to log in securely using their credentials. \\ \hline
	SRS-6 & The system shall use HTTPS/SSL for all data transfer between the client and server to prevent interception. \\ \hline
	SRS-7 & The system shall allow users to securely log out, terminating their current session. \\ \hline
\end{longtable}

\subsection{Category Management}

\subsubsection{Create Spending Category}
\begin{longtable}{|p{2.5cm}|p{12cm}|}
	\hline
	\textbf{ID} & \textbf{Requirement Description} \\
	\hline
	SRS-8 & The system shall allow users to create new, personalized spending category labels (e.g., "Groceries", "Utilities"). \\ \hline
	SRS-9 & The system shall validate that the category name is not a duplicate within the user's profile. \\ \hline
	SRS-10 & The system shall ensure all created categories are immediately available for the AI Sorting Engine to use. \\ \hline
\end{longtable}

\subsubsection{Manage Spending Categories}
\begin{longtable}{|p{2.5cm}|p{12cm}|}
	\hline
	\textbf{ID} & \textbf{Requirement Description} \\
	\hline
	SRS-11 & The system shall allow users to view a list of all their existing categories. \\ \hline
	SRS-12 & The system shall allow users to edit the name of an existing category. \\ \hline
	SRS-13 & The system shall allow users to delete a category (and optionally reassign its associated transactions). \\ \hline
\end{longtable}

\subsection{Data Ingestion Management}

\subsubsection{Upload Transaction File (CSV)}
\begin{longtable}{|p{2.5cm}|p{12cm}|}
	\hline
	\textbf{ID} & \textbf{Requirement Description} \\
	\hline
	SRS-14 & The system shall provide a simple interface for the user to upload a transaction history file. \\ \hline
	SRS-15 & The system is constrained to accepting only CSV (Comma-Separated Values) file formats. \\ \hline
	SRS-16 & The system shall provide a progress indicator or status feedback during the upload process. \\ \hline
\end{longtable}

\subsubsection{File Validation}
\begin{longtable}{|p{2.5cm}|p{12cm}|}
	\hline
	\textbf{ID} & \textbf{Requirement Description} \\
	\hline
	SRS-17 & The system shall check the uploaded file for necessary columns (specifically: Date, Description, and Amount). \\ \hline
	SRS-18 & The system shall provide clear, user-friendly error messages if the file structure is invalid or missing data. \\ \hline
	SRS-19 & The system shall reject files that exceed a predefined size limit to maintain performance. \\ \hline
\end{longtable}

\subsection{Intelligent Sorting Management}

\subsubsection{Process Transactions (AI Sorting)}
\begin{longtable}{|p{2.5cm}|p{12cm}|}
	\hline
	\textbf{ID} & \textbf{Requirement Description} \\
	\hline
	SRS-20 & The system shall use a smart classification model (AI) to read transaction descriptions from the uploaded file. \\ \hline
	SRS-21 & The system shall automatically assign the most relevant custom category to each transaction. \\ \hline
	SRS-22 & The system shall utilize Gemini API for the classification logic. \\ \hline
\end{longtable}

\subsubsection{Data Persistence \& Saving}
\begin{longtable}{|p{2.5cm}|p{12cm}|}
	\hline
	\textbf{ID} & \textbf{Requirement Description} \\
	\hline
	SRS-23 & The system shall securely save the processed transaction data, including the new category label, to the user's database. \\ \hline
	SRS-24 & The system shall ensure that once categorized, the financial history persists until explicitly deleted by the user. \\ \hline
\end{longtable}

\subsection{Analytics Management}

\subsubsection{View Dashboard}
\begin{longtable}{|p{2.5cm}|p{12cm}|}
	\hline
	\textbf{ID} & \textbf{Requirement Description} \\
	\hline
	SRS-25 & The system shall display a primary BI (Business Intelligence) dashboard view after processing is complete. \\ \hline
	SRS-26 & The dashboard shall be interactive, allowing users to hover over elements for more details. \\ \hline
\end{longtable}

\subsubsection{View Spending Breakdown}
\begin{longtable}{|p{2.5cm}|p{12cm}|}
	\hline
	\textbf{ID} & \textbf{Requirement Description} \\
	\hline
	SRS-27 & The system shall display clear charts (e.g., pie charts) visualizing total spending broken down by custom categories. \\ \hline
	SRS-28 & The system shall calculate and display total expenditure for the uploaded period. \\ \hline
\end{longtable}

\subsection{Administrator Management}

\subsubsection{Admin Panel Access}
\begin{longtable}{|p{2.5cm}|p{12cm}|}
	\hline
	\textbf{ID} & \textbf{Requirement Description} \\
	\hline
	SRS-29 & The System Administrator must use distinct credentials to log in to the dedicated administrative interface. \\ \hline
	SRS-30 & The Admin Panel shall be secured and inaccessible to standard users. \\ \hline
\end{longtable}

\subsubsection{User Account Management}
\begin{longtable}{|p{2.5cm}|p{12cm}|}
	\hline
	\textbf{ID} & \textbf{Requirement Description} \\
	\hline
	SRS-31 & The Admin shall have the ability to View and Search for any registered user account. \\ \hline
	SRS-32 & The Admin shall have the ability to Create, Edit, or Delete user accounts if necessary. \\ \hline
	SRS-33 & The Admin shall have read-only access to view user-created Categories and processed Transactions for auditing purposes. \\ \hline
\end{longtable}

\section{Non-Functional Requirements}

\subsection{Security}
\begin{itemize}
	\item \textbf{Account Protection:} User passwords must be stored using a strong one-way hash.
	\item \textbf{Session Security:} The system must use HTTPS for all communication to prevent data interception.
	\item \textbf{User Isolation:} A user's data must be completely isolated and inaccessible to any other user.
\end{itemize}

\subsection{Usability}
\begin{itemize}
	\item The interface must be intuitive, clean, and fully mobile-responsive.
	\item Training time for a normal user to understand the upload and sort process should be minimal (intuitive design).
\end{itemize}

\subsection{Reliability}
\begin{itemize}
	\item The system must be available and function correctly 99.5\% of the time.
	\item In any error scenario, the system must provide an appropriate message and retain the state of the user's data (no silent failures).
\end{itemize}

\subsection{Performance}
\begin{itemize}
	\item \textbf{File Processing Time:} A standard file (up to 500 transactions) must be processed, sorted by the AI, and ready for viewing within 60 seconds.
	\item The system dashboard should load visualization results in under 3 seconds after processing is complete.
\end{itemize}

\subsection{Design Constraints}
\begin{itemize}
	\item \textbf{Technology Stack:} Must use Python (Django Framework) for the backend logic.
	\item \textbf{Database:} Designed for use with PostgreSQL for scalability.
	\item \textbf{AI Service:} The system must utilize Gemini API for intelligent transaction classification.
	\item \textbf{Architecture:} The code must follow the MVT (Model-View-Template) architectural pattern.
\end{itemize}

\subsection{Data Ownership \& Business Rules}
\begin{itemize}
	\item \textbf{Data Ownership:} All uploaded data remains the intellectual property of the User.
	\item \textbf{No Third-Party Sharing:} User financial data may not be shared, sold, or exposed to any third party.
	\item \textbf{Persistence:} Processed data remains available until the user explicitly decides to delete it.
\end{itemize}

\section{External Interface Requirements}

\subsection{User Interfaces}
The user interface will be entirely web-based, optimized for responsive viewing on desktop and mobile devices. Key interfaces include:
\begin{itemize}
	\item Login/Registration Forms.
	\item Custom Category Management Form (CRUD interface).
	\item CSV File Upload Interface with progress/status feedback.
	\item Interactive BI Dashboard embedded through Metabase.
	\item A secure Admin Panel (for System Administrator use only).
\end{itemize}

\subsection{Hardware Interfaces}
No specialized hardware is required. The system will interface with standard internet-supported client devices (PC, laptop, mobile, tablet).

\subsection{Software Interfaces}
The system interfaces with the following software components:
\begin{itemize}
	\item \textbf{Python / Django:} Primary backend web application and API interface.
	\item \textbf{PostgreSQL:} Database management system for data persistence.
	\item \textbf{Python CSV and Django Forms:} Used for data parsing, validation, and persistence into database models.
	\item \textbf{Gemini API:} The core AI service used for transaction classification.
	\item \textbf{Metabase:} Business Intelligence tool used for dashboard visualization and embedded analytics.
\end{itemize}

\section{Detailed Functional Requirement Scenarios and Edge Cases}
To ensure the system remains stable under varied real-world usage, the core functional requirements must be evaluated against realistic data flows, alternative execution paths, and exception boundaries. This section outlines the sequential operations, alternative states, and failure recovery protocols for the four primary processing hubs of the application.

\subsection{Sign Up and Authentication Scenarios}
The user onboarding process represents the initial security gate of the application. To protect personal financial data, registration and authentication must enforce strict boundaries at the database and application levels.
\subsubsection{Sequential Processing Logic}
The registration sequence is initiated when a visitor requests the signup form. The Django backend generates a unique Cross-Site Request Forgery (CSRF) token, binds it to the session, and renders the form fields. The user must provide a username, a unique email address, a password, and a password confirmation. Upon submission, the backend performs form-level checks:
\begin{enumerate}
	\item \textbf{Syntax Verification:} The email address is checked against a standard format regex pattern to confirm structural validity.
	\item \textbf{Duplicate Auditing:} The database is queried to ensure the email address does not match any existing record.
	\item \textbf{Complexity Evaluation:} The password string is audited against complexity rules, verifying it contains at least eight characters, including numerical and alphabetical values.
	\item \textbf{Password Matching:} The password and confirmation strings are checked for character-for-character equality.
\end{enumerate}
If all validations pass, the Django authentication module hashes the password string using the PBKDF2 algorithm with a SHA-256 hash. The system opens a database transaction, saves the user record, generates a personalized default category set (e.g., Utilities, Groceries, Transport), commits the transaction, and redirects the user to the login screen.

\subsubsection{Alternate Courses and Exception Handling}
During registration, input errors or system faults can redirect the execution flow:
\begin{itemize}
	\item \textbf{Duplicate Account Conflict (Alternate Course):} If the database query identifies an existing user account registered under the submitted email, the system halts the creation process. The backend raises a form validation error, re-renders the registration page, and displays a notice stating that the email address is already in use.
	\item \textbf{Password Complexity Failure (Alternate Course):} If the password contains fewer than eight characters or fails complexity rules, the validation script blocks database entry. The system flags the password field and displays a message explaining the length and character requirements.
	\item \textbf{Database Transaction Timeout (Exception Path):} In the event of high database lock contention or server latency, the registration write operation may time out. The application catches the database transaction error, rolls back any partial database updates, logs the system error, and redirects the user to a page explaining that the server is temporarily busy, prompting them to try again.
\end{itemize}

\subsection{CSV Ingestion and Parsing Scenarios}
The ingestion pipeline converts external financial statement files into internal database records. Since banks and wallet services generate statement CSVs with different layouts, the ingestion service must handle files with structural variations.
\subsubsection{Sequential Processing Logic}
When an authenticated user requests a file upload, the application validates the file extension, selection of ingestion pathway, and size before opening the data stream. The ingestion pipeline supports uploading multiple CSV files simultaneously. If the files pass the initial boundaries, the Pandas library is initialized to read the CSV content of each file. The sequence executes as follows:
\begin{enumerate}
	\item \textbf{File Stream Buffering:} The uploaded files are read into memory as byte streams, avoiding disk-write latency.
	\item \textbf{Pathway Routing and Header Extraction:} The system routes each file stream based on the selected import pathway (Generic CSV Template or NayaPay Mobile Statement). For standard files, the parser matches columns against expected fields (Date, Description, and Amount) using case-insensitive mapping and predefined aliases. For NayaPay statements, the parser automatically maps specialized wallet export fields.
	\item \textbf{Data Cleaning:} The system loops through each row in the dataframe. It trims trailing spaces from descriptions, extracts transaction references, and removes currency symbols (such as Rs. or PKR) and thousands-separator commas from the amount values.
	\item \textbf{Transaction Batching:} The cleaned records are converted into transaction model instances and appended to a list for bulk database insertion.
\end{enumerate}
Once the batch list is fully populated, the Django ORM executes a single bulk insert operation to save the transactions.

\subsubsection{Alternate Courses and Exception Handling}
Data parsing must account for incomplete files, irregular cell contents, and format mismatches:
\begin{itemize}
	\item \textbf{Missing Column Headers (Alternate Course):} If the header row lacks one of the required columns (Date, Description, or Amount) and no alias matches are found, the Pandas parser halts. The system aborts the upload, discards the buffered stream, and renders a message stating which mandatory column is missing.
	\item \textbf{Non-Numeric Amount Formatting (Alternate Course):} When a row contains non-numeric text in the Amount column (such as ``N/A'' or ``FREE''), the numeric coercion module converts the entry to a null value. The row filter detects this null value and skips the row, recording it in the import log while continuing to process the other rows.
	\item \textbf{Encoding Mismatch Exception (Exception Path):} Financial statements may be encoded in UTF-8 or legacy formats like ISO-8859-1. If the parser encounters characters it cannot read, a decoding exception is raised. The application catches the exception, closes the stream, and displays a message asking the user to save the CSV in standard UTF-8 format.
\end{itemize}

\subsection{AI Classification and Inference Scenarios}
The Gemini API classification engine automates the assignment of spending categories using semantic natural language processing.
\subsubsection{Sequential Processing Logic}
When the user triggers the AI sorting function, the backend verifies that the user has custom categories and that there are uncategorized transactions. The classification sequence proceeds as follows:
\begin{enumerate}
	\item \textbf{Taxonomy Loading:} The system fetches the list of custom categories defined by the user.
	\item \textbf{Transaction Batching:} Uncategorized transactions are divided into batches of 25 to balance API payload size and network response latency.
	\item \textbf{Prompt Construction:} The engine compiles a structured prompt containing the candidate labels and the batch of transactions.
	\item \textbf{Gemini Classification:} The text descriptions are evaluated against the user's category list. The Gemini model calculates a confidence score for each category.
	\item \textbf{Confidence Assessment:} The engine checks the highest confidence score against the 0.5 safety threshold. If it exceeds the threshold, the category is assigned.
\end{enumerate}
The database is then updated with the category assignments.

\subsubsection{Alternate Courses and Exception Handling}
The AI sorting pipeline must handle low-confidence results and service interruptions:
\begin{itemize}
	\item \textbf{Low-Confidence Bypass (Alternate Course):} If the highest confidence score for a transaction falls below 0.5, the engine flags the prediction as uncertain. The system assigns the default label ``Uncategorized'' and stores the model's top guess as a suggestion, allowing the user to review the record later.
	\item \textbf{Empty Taxonomy Interception (Alternate Course):} If a user runs the classification engine without defining any categories, the application bypasses model execution. It marks all pending rows as ``Uncategorized'' and notifies the user to configure custom categories.
	\item \textbf{API Request Timeout (Exception Path):} High network latency or API rate limits may cause classification requests to time out. The application is configured to retry the request once. If the retry fails, the engine halts, rolls back the batch update, and displays a warning to the user, ensuring that no partially sorted records are left in the database.
\end{itemize}

\subsection{Dashboard Aggregation and Visualization Scenarios}
The analytics module provides secure data visualization dashboards using embedded Metabase resources.
\subsubsection{Sequential Processing Logic}
When a user accesses the dashboard, the application generates a secure, token-authorized iframe. The sequence runs as follows:
\begin{enumerate}
	\item \textbf{Configuration Audit:} The system checks the environment configuration settings to retrieve the Metabase site URL and embedding secret key.
	\item \textbf{Token Signature:} The backend compiles a JSON Web Token signed with the HMAC-SHA256 algorithm. The token contains the active user's identification number as a query filter parameter.
	\item \textbf{URL Construction:} The system appends the signed token to the Metabase dashboard site path, creating a secure embedding URL.
	\item \textbf{Frontend Embedding:} The Django template receives the URL and renders the Metabase dashboard inside an iframe within a glassmorphic dashboard container. The Metabase server filters the database queries dynamically using the user ID parameter.
\end{enumerate}

\subsection{Alternate Courses and Exception Handling}
The dashboard aggregation system must handle empty datasets and authorization issues:
\begin{itemize}
	\item \textbf{Missing Key Interception (Alternate Course):} If the Metabase secret key is missing from environment configurations, the dashboard view intercepts the error. It blocks the iframe generation and renders a warning banner explaining how to configure the environment settings.
	\item \textbf{Expired Session Token (Alternate Course):} The JSON Web Token contains an expiration timestamp set to 10 minutes. If the user keeps the dashboard open past this limit, Metabase displays a signature expired error. Refreshing the dashboard page generates a fresh token, restoring the view.
	\item \textbf{Zero Record State (Exception Path):} If the user has no transactions, the database queries return empty sets. Metabase displays empty visual states on the dashboard, explaining that data needs to be uploaded to populate the charts.
\end{itemize}
